%% toc: P3 35 35.9
\section{Proofs: failures of compactness and pinned ordinals}
\label{sec:annex-p03-w03-p03-041-compactness-failures}

Throughout, signatures, fields and well-orderings are as in
\cref{def:p03-w03-p03-041-setting}. Inductions on formulas are inductions
on the well-founded formula trees of \(L_{\omega_1\omega}\). The infinitary
clauses are handled in the same way as the finite ones, because the
satisfaction clauses for \(\bigwedge\) and \(\bigvee\) quantify over the
members of the index set (\cref{def:p02-lomega-one-omega}). For an element
\(a\) of a structure with a binary relation \(<\), write
\(\operatorname{pred}(a)=\{b:b<a\}\). If \(W\) is well-ordered by \(<\)
and \(h\) is the isomorphism of \(W\) onto its order type, then for every
\(b\in W\) the set \(\{c\in W:c<b\}\) is mapped by \(h\) onto
\(\{\gamma:\gamma<h(b)\}=h(b)\), so it has order type \(h(b)\). This fact
is used repeatedly below as the \emph{position fact}.

\begin{annexitem}{annex:p03-w03-p03-041-relativization}{thm:p03-w03-p03-041-relativization}{Renaming, relativization and reducts}
\begin{proof}
\emph{Part 1.} Induct on \(\varphi\). For an atomic
\(S(v_{i_1},\ldots,v_{i_k})\), its renaming is
\(S(v_{2i_1},\ldots,v_{2i_k})\). It holds at \(s\) iff
\((s(v_{2i_1}),\ldots,s(v_{2i_k}))\in S^{\mathfrak A}\), that is, iff
\(((s\circ e)(v_{i_1}),\ldots,(s\circ e)(v_{i_k}))\in S^{\mathfrak A}\).
This is the condition for the original atom at \(s\circ e\). Equations
are treated in the same way. Negation, binary connectives, \(\bigwedge\)
and \(\bigvee\) commute with renaming, and their satisfaction clauses
refer only to the components, so the induction hypothesis carries them.
Now take \(\exists v_i\,\psi\), whose renaming is
\(\exists v_{2i}\,\psi^{e}\). Write \(s[v\mapsto b]\) for the assignment
that differs from \(s\) at most by giving \(v\) the value \(b\). Because
\(e\) is injective,
\((s[v_{2i}\mapsto b])\circ e=(s\circ e)[v_i\mapsto b]\). At \(v_i\) both
sides give \(b\). At \(v_j\) with \(j\neq i\) both give \(s(v_{2j})\),
since \(2j\neq2i\). Hence \(\mathfrak A\entails\exists v_{2i}\,\psi^{e}[s]\)
iff some \(b\) has \(\mathfrak A\entails\psi^{e}[s[v_{2i}\mapsto b]]\). By
the induction hypothesis this holds iff some \(b\) has
\(\mathfrak A\entails\psi[(s\circ e)[v_i\mapsto b]]\), that is, iff
\(\mathfrak A\entails\exists v_i\,\psi[s\circ e]\). The universal case is
the same with ``every \(b\)'' in place of ``some \(b\)''. For a sentence,
neither side depends on the assignment, so \(\varphi\) and
\(\varphi^{e}\) have the same models.

\emph{Part 2.} Induct on \(\varphi\), for all assignments at once. Fix
\(s\) with \(s(v_1)=a\) and \(s(v)\in F_a\) for every
\(v\in\FV(\varphi)\), and \(s'\) into \(F_a\) agreeing with \(s\) on
\(\FV(\varphi)\). For an atomic \(S(u_1,\ldots,u_k)\) with
\(S\in\tau\), every \(u_j\) is free in it, so
\(s(u_j)=s'(u_j)\in F_a\). Since
\(S^{\mathfrak A_a}=S^{\mathfrak A}\cap F_a^{k}\), the atom holds in
\(\mathfrak A\) at \(s\) iff it holds in \(\mathfrak A_a\) at \(s'\).
An equation \(u=u'\) is evaluated at the same two elements on both sides.
Relativization leaves atoms unchanged, so the atomic case is complete. For
negation, binary connectives, \(\bigwedge\Psi\) and \(\bigvee\Psi\), each
component has its free variables among those of the compound. The
hypotheses on \(s\) and \(s'\) therefore hold for each component, and the
induction hypothesis together with the matching satisfaction clauses on
both sides gives the claim.

For \(\exists y\,\psi\), the variable \(y\) is not \(v_1\), because
\(v_1\) does not occur in \(\varphi\). So \(s[y\mapsto b](v_1)=a\) for
every \(b\), and
\(\mathfrak A\entails\exists y\,(R(v_1,y)\land\psi^{R})[s]\) iff some
\(b\in F_a\) has \(\mathfrak A\entails\psi^{R}[s[y\mapsto b]]\). For such
\(b\), the assignment \(s[y\mapsto b]\) sends every variable of
\(\FV(\psi)\subseteq\FV(\exists y\,\psi)\cup\{y\}\) into \(F_a\) and
agrees there with \(s'[y\mapsto b]\). By the induction hypothesis the
condition holds iff some \(b\in F_a\) has
\(\mathfrak A_a\entails\psi[s'[y\mapsto b]]\), that is, iff
\(\mathfrak A_a\entails\exists y\,\psi[s']\). For \(\forall y\,\psi\),
\(\mathfrak A\entails\forall y\,(R(v_1,y)\to\psi^{R})[s]\) iff every
\(b\in F_a\) has \(\mathfrak A\entails\psi^{R}[s[y\mapsto b]]\). By the
same use of the induction hypothesis this holds iff every \(b\in F_a\) has
\(\mathfrak A_a\entails\psi[s'[y\mapsto b]]\), which is
\(\mathfrak A_a\entails\forall y\,\psi[s']\). The domain of
\(\mathfrak A_a\) is exactly \(F_a\), so this is the universal clause in
\(\mathfrak A_a\).

\emph{Part 3.} The two structures have the same domain and interpret every
symbol of \(\tau\) identically. An induction on \(\varphi\), whose atomic
case uses only symbols of \(\tau\), therefore gives the same truth value at
every assignment.
\end{proof}
\end{annexitem}

\begin{annexitem}{annex:p03-w03-p03-041-pinning-compactness}{thm:p03-w03-p03-041-pinning-compactness}{Countable compactness and pinned ordinals}
\begin{proof}
Write \(\operatorname{Fld}(x)\) for \(\exists y\,(x<y\lor y<x)\), and
\(\operatorname{LO}_{\mathrm f}\) for the first-order sentence saying
that \(<\) is irreflexive and transitive and that any two elements
satisfying \(\operatorname{Fld}\) are equal or comparable. Its models are
exactly the structures in which \(<\) is a strict linear order of its
field.

\emph{Step 0: \(w_\kappa(\mathcal F)\geq\omega\).} The first-order
sentence \(\forall x\,\forall y\,\neg x<y\) pins down \(0\). Its models
have empty field, and the empty relation well-orders the empty field with
type \(0\). For \(n\geq2\), let \(\rho_n\) be the conjunction of
\(\operatorname{LO}_{\mathrm f}\) with a first-order sentence saying that
exactly \(n\) elements satisfy \(\operatorname{Fld}\). In a model of
\(\rho_n\), \(<\) is a strict linear order of an \(n\)-element field.
Every nonempty subset of a finite linear order has a least element, so
\(<\) is a well-ordering of type \(n\). The structure \((n,<)\) is a model
with field \(n\), since \(n\geq2\). By (F1) these are
\(\mathcal F\)-sentences, and each set \(\{\rho_n\}\) has cardinality
\(1\leq\kappa\). The supremum of \(\{0\}\cup\{n:n\geq2\}\) is \(\omega\).

\emph{Step A: from \(\alpha\) to \(\alpha+1\).} Let \(\Phi\) be a set of
\(\mathcal F\)-sentences of \(\tau\) that pins down \(\alpha\geq\omega\)
via \(<\). Let \(\prec\) be a new binary symbol, write
\(\operatorname{Least}(x)\) for
\(\operatorname{Fld}(x)\land\neg\exists y\,y<x\) and
\(\operatorname{NL}(x)\) for
\(\operatorname{Fld}(x)\land\neg\operatorname{Least}(x)\), and let
\(\chi\) be
\[
 \forall x\,\forall y\Bigl(x\prec y\leftrightarrow
   \operatorname{NL}(x)\land\bigl((\operatorname{NL}(y)\land x<y)
   \lor\operatorname{Least}(y)\bigr)\Bigr).
\]
This sentence is first-order, so \(\Phi\cup\{\chi\}\) is a set of
\(\mathcal F\)-sentences of \(\tau\cup\{\prec\}\) by (F1).

Let \(\mathfrak A\) be a model of \(\Phi\cup\{\chi\}\). By
\cref{thm:p03-w03-p03-041-relativization}, part~3, its reduct satisfies
\(\Phi\), so \(<^{\mathfrak A}\) well-orders its field \(X\). If \(X\) is
empty, then \(\prec^{\mathfrak A}\) is empty and well-orders its empty
field. Otherwise \(X\) has a least element \(l\) and at least two elements,
and \(\chi\) makes \(x\prec y\) exactly in two cases: \(x,y\in X\setminus\{l\}\)
with \(x<y\), or \(x\in X\setminus\{l\}\) and \(y=l\). Every element of
\(X\) occurs in such a pair, so the field of \(\prec\) is \(X\).
Irreflexivity holds because \(x\prec x\) would need \(x\) to be
non-least and also either \(x<x\) or \(x=l\). The element \(l\) never
stands on the left of \(\prec\). Totality holds because two distinct
non-least elements are \(<\)-comparable and every non-least element is
\(\prec l\). For transitivity, let \(x\prec y\prec z\). Then \(x\) and
\(y\) are non-least. If \(z=l\), then \(x\prec z\). If \(z\neq l\), then
\(x<y<z\), so \(x<z\) and \(x\prec z\). Now let \(Y\subseteq X\) be nonempty.
If \(Y\setminus\{l\}\) is nonempty, its \(<\)-least element is
\(\prec\)-least in \(Y\), because it lies below \(l\). If not, then
\(Y=\{l\}\). So \(\prec^{\mathfrak A}\) well-orders its field.

Take a model \(\mathfrak A\) of \(\Phi\) whose field has type \(\alpha\),
with isomorphism \(f\) onto \(\alpha\), and interpret \(\prec\) as \(\chi\)
requires. The map \(g\) on \(X\setminus\{l\}\) with \(g(c)=f(c)-1\) when
\(f(c)\) is finite and \(g(c)=f(c)\) when \(f(c)\geq\omega\) is an order
isomorphism onto \(\alpha\). It is order preserving, and it is onto
because \(\alpha\geq\omega\). Extending \(g\) by \(l\mapsto\alpha\) gives
an isomorphism of \((X,\prec)\) onto \(\alpha+1\). So
\(\Phi\cup\{\chi\}\) pins down \(\alpha+1\), and it is countable when
\(\Phi\) is.

\emph{Step B: an infinite pinned ordinal defeats compactness.} Let
\(\Phi\) be a set of \(\mathcal F\)-sentences of \(\tau\) that pins down
some \(\alpha\geq\omega\) via \(<\). Add new unary symbols \(U_n\)
(\(n<\omega\)), and let \(C\) consist of the first-order sentences
\[
 \operatorname{Sg}_n:\ \exists x\,\forall y\,(U_n(y)\leftrightarrow y=x),
 \qquad
 \operatorname{Dn}_n:\ \exists x\,\exists y\,(U_{n+1}(x)\land U_n(y)\land x<y).
\]
By (F1), \(\Phi\cup C\) is a set of \(\mathcal F\)-sentences. We show that
every finite subset has a model but the whole set has none.

Let \(\Psi_0\subseteq\Phi\cup C\) be finite, and choose \(N<\omega\) larger
than every \(n\) for which \(\operatorname{Sg}_n\) or
\(\operatorname{Dn}_n\) lies in \(\Psi_0\). Take a model \(\mathfrak A\)
of \(\Phi\) whose field has type \(\alpha\), with isomorphism \(f\) onto
\(\alpha\). Since \(N<\omega\leq\alpha\), the elements
\(a_i=f^{-1}(N-i)\) for \(i\leq N\) exist and satisfy
\(a_N<a_{N-1}<\cdots<a_0\). Expand \(\mathfrak A\) by
\(U_n=\{a_{\min(n,N)}\}\). Every \(\operatorname{Sg}_n\) holds. Each
\(\operatorname{Dn}_n\) in \(\Psi_0\) has \(n+1\leq N\) and holds because
\(a_{n+1}<a_n\). The members of \(\Phi\) hold by part~3 of
\cref{thm:p03-w03-p03-041-relativization}.

Suppose \(\mathfrak B\) were a model of \(\Phi\cup C\). Let \(u_n\) be the
unique element of \(U_n^{\mathfrak B}\). Then \(u_{n+1}<u_n\) for every
\(n\). Each \(u_n\) is in the field, and \(\{u_n:n<\omega\}\) is a
nonempty subset of the field with no least element. But the reduct of
\(\mathfrak B\) to \(\tau\) is a model of \(\Phi\) by part~3 of
\cref{thm:p03-w03-p03-041-relativization}, and it has the same \(<\). So
by (i) of \cref{def:p03-w03-p03-041-pinning} this \(<\) is a
well-ordering of its field, which is a contradiction. This step uses no
form of choice.

\emph{\((1)\Rightarrow(2)\).} If a countable \(\Phi\) pinned down an
infinite ordinal, then \(\Phi\cup C\) would be a countable set violating
\(\aleph_0\)-compactness, by Step~B.

\emph{\((2)\Rightarrow(3)\).} Under (2), every ordinal pinned down by a
countable set is finite, so \(w_{\aleph_0}(\mathcal F)\leq\omega\). Step~0
gives equality.

\emph{\((3)\Rightarrow(2)\).} If a countable set pinned down an infinite
\(\alpha\), then by Step~A a countable set would pin down \(\alpha+1\), and
so \(w_{\aleph_0}(\mathcal F)\geq\alpha+1>\omega\).

\emph{\((2)\Rightarrow(1)\), and the set that pins down \(\omega\).}
Suppose \(\mathcal F\) is not \(\aleph_0\)-compact. Take a positive
relational \(\tau\) and a countable set \(\Phi\) of
\(\mathcal F\)-sentences of \(\tau\) in which every finite subset has a
model but \(\Phi\) has none. Then \(\Phi\) is nonempty, because any
\(\tau\)-structure is a model of the empty set. List it as
\(\Phi=\{\varphi_n:n<\omega\}\), with repetitions allowed. Let \(<\) and
\(R\) be new binary symbols, \(\tau'=\tau\cup\{<,R\}\), and
\(\psi_n=(\varphi_n^{e})^{R}\). This is defined because \(v_1\) does not
occur in \(\varphi_n^{e}\), and its only free variable is \(v_1\). For
\(n\geq1\), let \(\operatorname{Le}_n(v_1)\) be the first-order formula
\[
 \exists y_1\cdots\exists y_n\Bigl(\bigwedge_{1\leq i<j\leq n}\neg y_i=y_j
   \land\bigwedge_{i=1}^{n}(y_i<v_1\lor y_i=v_1)\Bigr),
\]
with \(y_1,\ldots,y_n\) distinct from \(v_1\). It says that at least \(n\)
elements are \(\leq v_1\). Let \(\operatorname{Le}_0(v_1)\) be
\(v_1=v_1\). Let \(\Phi'\) consist of
\[
\begin{split}
 &(1)\ \operatorname{LO}_{\mathrm f},\qquad
  (2)\ \forall v_1\bigl(\operatorname{Fld}(v_1)\to\exists y\,R(v_1,y)\bigr),\\
 &(3_n)\ \forall v_1\bigl(\operatorname{Fld}(v_1)\land
        \operatorname{Le}_n(v_1)\to\psi_n\bigr)\quad(n<\omega).
\end{split}
\]
By (F3), \(\varphi_n^{e}\in\mathcal F[\tau]\). By (F4),
\(\psi_n\in\mathcal F[\tau\cup\{R\}]\subseteq\mathcal F[\tau']\), using
(F1). Writing \(\theta\to\psi\) as \(\neg(\theta\land\neg\psi)\) and
\(\forall\) as \(\neg\exists\neg\), (F1) and (F2) put every member of
\(\Phi'\) in \(\mathcal F[\tau']\). The set \(\Phi'\) is countable.

\emph{Every model of \(\Phi'\) well-orders its field, with type at most
\(\omega\).} Let \(\mathfrak A\entails\Phi'\), and suppose some field
element \(a\) had infinitely many elements \(\leq a\). Then
\(\mathfrak A\entails\operatorname{Le}_n[a]\) for every \(n\), so (3\(_n\))
gives \(\mathfrak A\entails\psi_n[s]\) for every \(n\), where
\(s(v_1)=a\). By (2), \(F_a\) is nonempty. Part~2 of
\cref{thm:p03-w03-p03-041-relativization}, applied to the sentence
\(\varphi_n^{e}\), gives \(\mathfrak A_a\entails\varphi_n^{e}\). Part~1
then gives \(\mathfrak A_a\entails\varphi_n\). So the \(\tau\)-structure
\(\mathfrak A_a\) is a model of \(\Phi\), contrary to the choice of
\(\Phi\). Hence every field element has only finitely many elements below
it. Let \(Y\) be a nonempty subset of the field, and take \(b\in Y\). The
finitely many members of \(Y\) that are \(\leq b\) form a finite linear
order, and its least element is least in \(Y\). So \(<^{\mathfrak A}\)
well-orders its field. The map sending \(a\) to the number of elements
below it is the isomorphism of the field onto its type, and its values are
natural numbers. So the type is at most \(\omega\).

\emph{A model of type \(\omega\).} For each \(k<\omega\), the finite set
\(\{\varphi_0,\ldots,\varphi_{k+1}\}\) has a model. Using the Axiom of
Choice, choose one, \(\mathfrak B_k\). Let the domain be
\[
 A=(\omega\times\{0\})\cup\bigcup_{k<\omega}(B_k\times\{k+1\}),
\]
with \((j,0)<(k,0)\) iff \(j<k\), and with
\(R=\{((k,0),(b,k+1)):k<\omega,\ b\in B_k\}\). For each \(S\in\tau\) of
arity \(m\geq1\), let \(S^{\mathfrak A}\) be the set of tuples
\(((b_1,k+1),\ldots,(b_m,k+1))\) with
\((b_1,\ldots,b_m)\in S^{\mathfrak B_k}\). The field of \(<\) is
\(\omega\times\{0\}\), ordered with type \(\omega\), so (1) holds. The
\(R\)-successors of \((k,0)\) form \(B_k\times\{k+1\}\), which is nonempty,
so (2) holds. Every tuple in \(S^{\mathfrak A}\) has all its entries in a
single set \(B_k\times\{k+1\}\), because \(m\geq1\) and each tuple carries
one tag. Hence the restriction of \(S^{\mathfrak A}\) to
\(B_k\times\{k+1\}\) is the tagged copy of \(S^{\mathfrak B_k}\). The map
\(b\mapsto(b,k+1)\) is therefore an isomorphism of \(\mathfrak B_k\) onto
\(\mathfrak A_{(k,0)}\).

An isomorphism preserves the satisfaction of every
\(L_{\omega_1\omega}\)-formula under corresponding assignments. This is
proved by induction on formulas. Atoms are preserved because the map
preserves and reflects every relation and is injective. The connective
and infinitary clauses are preserved componentwise. The quantifier clauses
are preserved because the map is a bijection. Now let \((k,0)\) satisfy
\(\operatorname{Le}_n\). There are \(k+1\) elements \(\leq(k,0)\), so
\(n\leq k+1\). Then \(\mathfrak B_k\entails\varphi_n\), hence
\(\mathfrak A_{(k,0)}\entails\varphi_n\), and so
\(\mathfrak A_{(k,0)}\entails\varphi_n^{e}\) by part~1. Part~2 gives
\(\mathfrak A\entails\psi_n[s]\) with \(s(v_1)=(k,0)\). So (3\(_n\))
holds. Hence \(\Phi'\) pins down \(\omega\), and (2) fails. This proves
\((2)\Rightarrow(1)\) by contraposition, together with the stated clause
about \(\omega\).

\emph{First-order logic.} Let \(\Phi\) be any set of first-order sentences
that pins down an ordinal \(\alpha\geq\omega\). By Step~B,
\(\Phi\cup C\) is a set of first-order sentences whose every finite subset
has a model, and it has none. This contradicts
\cref{thm:p03-compactness-ultraproduct}. So every ordinal pinned down by any set of
first-order sentences is finite, \(w_\kappa\leq\omega\) for every
\(\kappa\), and Step~0 gives equality.
\end{proof}
\end{annexitem}

\begin{annexitem}{annex:p03-w03-p03-041-mu}{thm:p03-w03-p03-041-mu}{What \(\mu_\beta\) says in a linear order}
\begin{proof}
Induct on \(\beta\). Assume the claim for every \(\gamma<\beta\), and let
\[
 S=\{b:\operatorname{pred}(b)\text{ is well-ordered with type }<\beta\}.
\]
By the induction hypothesis and the clause for \(\bigvee\),
\(\mathfrak A\entails\bigvee\{\mu_\gamma(y)\mid\gamma<\beta\}\) at \(b\)
iff \(b\in S\). Hence \(\mathfrak A\entails\mu_\beta[a]\) iff
\(\operatorname{pred}(a)=S\). For \(\beta=0\), the set \(S\) is empty and
the claim reads: \(\operatorname{pred}(a)=\varnothing\) iff
\(\operatorname{pred}(a)\) is well-ordered of type \(0\). This is true.
The argument below covers this case as well.

\emph{If \(\operatorname{pred}(a)\) is well-ordered with type \(\beta\),
then \(\operatorname{pred}(a)=S\).} Let \(h\) be the isomorphism of
\(\operatorname{pred}(a)\) onto \(\beta\). For \(b<a\), transitivity gives
\(\operatorname{pred}(b)=\{c\in\operatorname{pred}(a):c<b\}\). By the
position fact this set has type \(h(b)<\beta\), so \(b\in S\). Conversely,
let \(b\in S\). If \(b=a\), then \(\operatorname{pred}(b)\) has type
\(\beta\), which is not below \(\beta\). If \(a<b\), then
\(\operatorname{pred}(a)=\{c\in\operatorname{pred}(b):c<a\}\), and by the
position fact in the well-order \(\operatorname{pred}(b)\) its type
\(\beta\) is below the type of \(\operatorname{pred}(b)\), which is itself
below \(\beta\). Both cases are impossible. By totality \(b<a\), so
\(S\subseteq\operatorname{pred}(a)\).

\emph{If \(\operatorname{pred}(a)=S\), then \(S\) is well-ordered with
type \(\beta\).} Let \(X\subseteq S\) be nonempty and \(b\in X\). If no
member of \(X\) is below \(b\), then \(b\) is least in \(X\) by totality.
Otherwise \(X\cap\operatorname{pred}(b)\) is a nonempty subset of the
well-order \(\operatorname{pred}(b)\), since \(b\in S\). Its least element
\(c\) is least in \(X\): a member \(d\) of \(X\) is either below \(b\),
and then \(c\leq d\), or satisfies \(d\geq b>c\). So \(S\) is
well-ordered. Let \(h\) be its isomorphism onto its type \(\delta\). For
\(b\in S\), every \(c<b\) satisfies \(c<a\), so
\(\operatorname{pred}(b)=\{c\in S:c<b\}\). By the position fact this set
has type \(h(b)\), and \(h(b)<\beta\) because \(b\in S\). Since \(h\) is
onto \(\delta\), every \(\gamma<\delta\) is below \(\beta\), so
\(\delta\leq\beta\). If \(\delta<\beta\), then \(\operatorname{pred}(a)=S\)
is well-ordered with type below \(\beta\). That would put \(a\) in
\(S=\operatorname{pred}(a)\), giving \(a<a\), which contradicts
irreflexivity. Hence \(\delta=\beta\).
\end{proof}
\end{annexitem}

\begin{annexitem}{annex:p03-w03-p03-041-countable-ordinals}{thm:p03-w03-p03-041-countable-ordinals}{\(L_{\omega_1\omega}\) pins down every countable ordinal from \(2\) on}
\begin{proof}
Let \(\mathfrak A=(A,<)\entails\sigma_\alpha\). Then \(<\) is a strict
linear order of \(A\). By \cref{thm:p03-w03-p03-041-mu}, every \(a\in A\)
has \(\operatorname{pred}(a)\) well-ordered with some type below
\(\alpha\). The argument of the second half of the proof of
\cref{thm:p03-w03-p03-041-mu} applies, with \(A\) in place of \(S\), and
shows that \(A\) is well-ordered. Let \(h\) be the isomorphism of \(A\)
onto its type \(\delta\). By the position fact,
\(h(a)\) is the type of \(\operatorname{pred}(a)\), so \(h(a)<\alpha\)
for every \(a\), and \(h\) is onto \(\delta\), so \(\delta\leq\alpha\). Also
\(\delta\geq1\), because \(A\) is nonempty. Conversely, let \((A,<)\) be a
well-order of type \(\delta\) with \(1\leq\delta\leq\alpha\), with
isomorphism \(h\) onto \(\delta\). It satisfies
\(\operatorname{LO}(<)\). Each \(a\) has \(\operatorname{pred}(a)\) of
type \(h(a)<\delta\leq\alpha\), so by \cref{thm:p03-w03-p03-041-mu} the
disjunct \(\mu_{h(a)}\) holds at \(a\). So \((A,<)\entails\sigma_\alpha\).

Now let \(2\leq\alpha<\omega_1\). Every model of \(\sigma_\alpha\) is a
well-order of its whole domain. If the domain has at least two elements,
the field is the domain. If the domain has one element, the field is empty
and is well-ordered. So (i) of \cref{def:p03-w03-p03-041-pinning} holds.
The structure \((\alpha,\in)\) is a model, and since \(\alpha\geq2\) its
field is \(\alpha\), of type \(\alpha\). So (ii) holds. Since
\(\{\sigma_\alpha\}\) has one member,
\(w(L_{\omega_1\omega})=w_1(L_{\omega_1\omega})\geq
\sup\{\alpha:2\leq\alpha<\omega_1\}=\omega_1\).

Finally, \(\{\sigma_\omega\}\) is a countable set of
\(L_{\omega_1\omega}\)-sentences that pins down the infinite ordinal
\(\omega\). Step~B of the proof of
\cref{thm:p03-w03-p03-041-pinning-compactness} uses no choice. Applied to
\(\Phi=\{\sigma_\omega\}\), it produces a countable set of
\(L_{\omega_1\omega}\)-sentences whose every finite subset has a model and
which has none. So \(L_{\omega_1\omega}\) is not \(\aleph_0\)-compact.
\end{proof}
\end{annexitem}

\begin{annexitem}{annex:p03-w03-p03-041-transfer}{thm:p03-w03-p03-041-transfer}{Three first-order transfers that fail for \(L_{\omega_1\omega}\)}
\begin{proof}
\emph{Clause 1.} Let \(A\) be finite with \(n\geq1\) elements
\(a_1,\ldots,a_n\). The disjunct for \(n\) holds, witnessed by
\(v_i\mapsto a_i\), since every \(y\) equals some \(a_i\). Conversely, if
\(\mathfrak A\entails\varphi_{\mathrm{fin}}\), some disjunct \(n\) holds.
So there are \(a_1,\ldots,a_n\) such that every element equals one of
them, and \(A\) has at most \(n\) elements. The models of
\(\varphi_{\mathrm{fin}}\) are therefore exactly the finite structures, of
every finite size, and none is infinite.

For the first-order contrast, let \(\lambda_k\) be the first-order sentence
saying that there are at least \(k\) distinct elements. Let \(T\) be a
first-order theory with arbitrarily large finite models. A finite subset
of \(T\cup\{\lambda_k:k\geq1\}\) mentions \(\lambda_k\) only for
\(k\leq K\), for some \(K\), and a model of \(T\) with at least \(K\)
elements satisfies it. By \cref{thm:p03-compactness-ultraproduct}, the
whole set has a model, and that model is infinite.

\emph{Clause 2.} A nonprincipal ultrafilter \(\mathcal U\) on \(\omega\)
contains every cofinite set. Suppose a cofinite \(C\) were not in
\(\mathcal U\). Then its finite complement would be in \(\mathcal U\),
and it would be nonempty because \(\varnothing\notin\mathcal U\). An
ultrafilter that contains a finite union of sets contains one of them, so
some singleton would be in \(\mathcal U\), and \(\mathcal U\) would be
principal.

Identify \(\mathfrak M_n\) with \(\{0,\ldots,n\}\), and for \(j<\omega\) let
\(f_j(n)=\min(j,n)\). For \(i<j\), the set
\(\{n:f_i(n)=f_j(n)\}=\{n:n\leq i\}\) is finite. Its complement is cofinite
and lies in \(\mathcal U\), so the finite set does not. Hence the classes of
\(f_0,f_1,f_2,\ldots\) are pairwise distinct, and the ultraproduct is
infinite. By clause~1 it does not satisfy \(\varphi_{\mathrm{fin}}\). On
the other hand,
\(\{n:\mathfrak M_n\entails\varphi_{\mathrm{fin}}\}=\omega\in\mathcal U\),
by clause~1. So the equivalence asserted by \L o\'s's theorem for
first-order sentences fails for the sentence \(\varphi_{\mathrm{fin}}\).

\emph{Clause 3.} If \(\mathfrak M\entails\forall x\bigvee_{n<\omega}x=d_n\),
then every element is some \(d_n^{\mathfrak M}\). So
\(n\mapsto d_n^{\mathfrak M}\) maps \(\omega\) onto the domain, and the
domain is countable. The structure with domain \(\omega\) and
\(d_n=n\) satisfies the sentence, since each \(x\in\omega\) is
\(d_x\), and it is infinite. The first-order contrast is
\cref{thm:p02-upward-ls-elementary-extension}. Apply it to an infinite model
\(\mathcal M\) of a first-order theory with
\(\kappa=\max(\abs{\Lang},\abs{\dom\mathcal M},\aleph_1)\). It gives an
elementary extension of cardinality \(\kappa\geq\aleph_1\). That extension
satisfies the same first-order sentences, so it is an uncountable model of
the theory.
\end{proof}
\end{annexitem}

\begin{annexitem}{annex:p03-w03-p03-041-undefinability}{thm:p03-w03-p03-041-undefinability}{Undefinability of well-ordering}
\begin{proof}
The construction is Marker's
\cite[proof after Corollary~4.28, pp.~52--53]{Marker2014InfinitaryNotes},
with the changes of \cref{def:p03-w03-p03-041-marker-gaps}. Two facts
about \(L_{\omega_1\omega}\) are used throughout, each proved by induction
on formulas. First, whether a structure satisfies a formula at an
assignment depends only on the interpretations of the symbols that occur
in the formula. Second, a formula has only countably many subformulas,
since the subformulas of \(\psi\) are \(\psi\) and the subformulas of its
countably many immediate constituents. In particular \(\varphi\) contains
only countably many symbols.

\emph{Step 1: reduction to a countable signature.} Let \(\tau_0\) be the
set of symbols occurring in \(\varphi\) together with \(<\). It is
countable. The \(\tau_0\)-reduct of each model in the hypothesis is still
a model of \(\varphi\), with the same relation \(<\) and the same map
\(e\). Suppose a \(\tau_0\)-structure \(\mathfrak N_0\entails\varphi\)
with countable domain and a map \(h\) as required have been found.
Interpret every symbol of \(\tau\setminus\tau_0\) on the domain of
\(\mathfrak N_0\) in any way, which is possible because domains are
nonempty. The result is a \(\tau\)-structure \(\mathfrak N\) with the same
countable domain. It satisfies \(\varphi\) by the first fact, and \(h\)
keeps its property. So from now on \(\tau=\tau_0\) is countable.

\emph{Step 2: the extended signature.} Add to \(\tau\) a countably
infinite set \(C\) of new constant symbols and a further set
\(D=\{d_q:q\in\mathbb Q\}\) of new constant symbols, pairwise distinct.
Let \(\tau^{+}=\tau\cup C\cup D\). It is countable, and \(C\) is an
infinite set of constant symbols of \(\tau^{+}\), as
\cref{def:p03-w03-p03-039-consistency-property} requires. Put
\(\Delta_{\mathbb Q}=\{d_q<d_r:q,r\in\mathbb Q,\ q<r\}\).

For a set \(s\) of \(\tau^{+}\)-sentences, write \(C(s)\) and \(D(s)\)
for the sets of constants of \(C\) and of \(D\) that occur in members of
\(s\). Call \(s\) \emph{small} if \(s\) is finite and \(C(s)\cup D(s)\) is
finite. A single sentence may mention infinitely many constants, which is
why the second condition is separate. For a small \(s\), list
\(D(s)=\{d_{q_1},\ldots,d_{q_m}\}\) with \(q_1<\cdots<q_m\), so that
\(m=0\) when \(D(s)\) is empty.

\emph{Step 3: realizations.} Let \(s\) be small and let \(\alpha\) be an
ordinal. An \emph{\(\alpha\)-realization} of \(s\) consists of
\begin{enumerate}
\item a \(\tau\)-structure \(\mathfrak M\entails\varphi\);
\item an ordinal \(\gamma\) and a map \(e\colon\gamma\to\dom\mathfrak M\)
      with \(e(\xi)<^{\mathfrak M}e(\eta)\) whenever \(\xi<\eta<\gamma\);
\item ordinals \(\beta_1<\cdots<\beta_m<\gamma\) with
      \(\alpha\leq\beta_1\), with \(\beta_i+\alpha\leq\beta_{i+1}\) for
      \(1\leq i<m\), and with \(\beta_m+\alpha\leq\gamma\); when \(m=0\)
      the requirement is \(\alpha\leq\gamma\);
\item a map \(g\colon C(s)\to\dom\mathfrak M\),
\end{enumerate}
such that the expansion \(\mathfrak M^{\ast}\) of \(\mathfrak M\) to
\(\tau\cup C(s)\cup D(s)\) interpreting each \(c\in C(s)\) as \(g(c)\) and
each \(d_{q_i}\) as \(e(\beta_i)\) satisfies every sentence of \(s\). All
sentences of \(s\) are sentences of \(\tau\cup C(s)\cup D(s)\), so this
makes sense. Let \(\mathcal S\) be the set of small \(s\) that have an
\(\alpha\)-realization for every \(\alpha<\omega_1\).

\emph{Monotonicity.} If \(\alpha'\leq\alpha\), every
\(\alpha\)-realization of \(s\) is an \(\alpha'\)-realization of \(s\).
Indeed \(\alpha'\leq\alpha\leq\beta_1\), and ordinal addition is weakly
increasing in its right argument, so
\(\beta_i+\alpha'\leq\beta_i+\alpha\leq\beta_{i+1}\) and
\(\beta_m+\alpha'\leq\gamma\).

\emph{Truth fact.} Let \(\alpha\geq1\) and take an \(\alpha\)-realization
of a small \(s\), with expansion \(\mathfrak M^{\ast}\). Then
\(\mathfrak M^{\ast}\) satisfies every member of
\[
 T(s)=s\cup\{\varphi\}\cup\{d_q<d_r\in\Delta_{\mathbb Q}:d_q,d_r\in D(s)\}.
\]
For \(s\) this is the definition. For \(\varphi\) it is the first fact,
since \(\mathfrak M\entails\varphi\) and \(\varphi\) contains no new
constant. If \(d_{q_i},d_{q_j}\in D(s)\) with \(q_i<q_j\), then \(i<j\),
so \(\beta_i<\beta_j\) and \(e(\beta_i)<^{\mathfrak M}e(\beta_j)\).

\emph{Adding a sentence.} Let \(s\in\mathcal S\) and let \(\chi\) be a
\(\tau^{+}\)-sentence whose constants from \(C\cup D\) all lie in
\(C(s)\cup D(s)\). Suppose \(\chi\) holds in the expansion of every
\(\alpha\)-realization of \(s\) with \(\alpha\geq1\). Then
\(s\cup\{\chi\}\in\mathcal S\). For \(\alpha\geq1\), an
\(\alpha\)-realization of \(s\) is one of \(s\cup\{\chi\}\), since the
constants and their interpretations are unchanged. For \(\alpha=0\), a
\(1\)-realization serves, by monotonicity.

\emph{Step 4: adding constants.} For a small \(s\) and a finite
\(K\subseteq C\cup D\), let \(s[K]=s\cup\{k=k:k\in K\}\). It is small, and
\(C(s[K])\cup D(s[K])=C(s)\cup D(s)\cup K\). We show that
\(s\in\mathcal S\) implies \(s[K]\in\mathcal S\). Constants of \(K\)
already in \(C(s)\cup D(s)\) change nothing, since \(k=k\) holds in every
expansion. By induction on the size of \(K\), it therefore suffices to add
one constant \(k\notin C(s)\cup D(s)\).

If \(k=c\in C\), take any \(\alpha\)-realization of \(s\) and extend
\(g\) by sending \(c\) to any element. The sentences of \(s\) do not
mention \(c\), so they stay true, and \(c=c\) is true.

If \(k=d_r\in D\), then \(r\) differs from every \(q_i\). Let
\(1\leq\alpha<\omega_1\) and put \(\alpha^{\ast}=\alpha+\alpha\), a
countable ordinal. Take an \(\alpha^{\ast}\)-realization of \(s\) with
\(\gamma\), \(e\), \(\beta_1,\ldots,\beta_m\) and \(g\). Choose a position
\(b\) for \(d_r\) as follows.
\begin{itemize}
\item If \(m=0\) or \(r<q_1\), let \(b=\alpha\). Then \(\alpha\leq b\) and
      \(b+\alpha=\alpha^{\ast}\), which is at most \(\beta_1\) when
      \(m\geq1\) and at most \(\gamma\) when \(m=0\).
\item If \(q_j<r<q_{j+1}\), let \(b=\beta_j+\alpha\). Then
      \(\beta_j+\alpha\leq b\) and, by associativity,
      \(b+\alpha=\beta_j+\alpha^{\ast}\leq\beta_{j+1}\).
\item If \(m\geq1\) and \(r>q_m\), let \(b=\beta_m+\alpha\). Then
      \(\beta_m+\alpha\leq b\) and
      \(b+\alpha=\beta_m+\alpha^{\ast}\leq\gamma\).
\end{itemize}
Insert \(b\) into the list of positions at the place of \(r\) among the
\(q_i\). Every gap between consecutive old positions, and the gap from the
last old position to \(\gamma\), was at least \(\alpha^{\ast}\geq\alpha\),
and the first old position is at least \(\alpha^{\ast}\geq\alpha\). The
new gaps were checked in the three cases. Since \(\alpha\geq1\), each
position is strictly below the next one and below \(\gamma\), so the new
list is strictly increasing and \(e(b)\) is defined. Interpret \(d_r\) as
\(e(b)\). The sentences of \(s\) do not mention \(d_r\), and \(d_r=d_r\)
is true. This is an \(\alpha\)-realization of \(s[\{d_r\}]\). For
\(\alpha=0\), a \(1\)-realization serves.

\emph{Step 5: the consistency property.} Write
\(\hat s=s\cup\{\varphi\}\cup\Delta_{\mathbb Q}\) and let
\[
 \Sigma=\bigl\{\mu:\ \mu\subseteq\hat s\text{ for some }s\in\mathcal S\bigr\}.
\]
Each member is a countable set of \(\tau^{+}\)-sentences. The empty set is
small, and it is in \(\mathcal S\): given \(\alpha<\omega_1\), the
hypothesis gives \(\mathfrak M\entails\varphi\) and \(e\) on
\(\gamma=\alpha\), and with \(m=0\) and \(g\) empty the requirement
\(\alpha\leq\gamma\) holds. So \(\{\varphi\}\cup\Delta_{\mathbb Q}\in
\Sigma\), and \(\Sigma\) is nonempty.

We verify the clauses of \cref{def:p03-w03-p03-039-consistency-property}.
Fix \(\mu\in\Sigma\) and \(s\in\mathcal S\) with \(\mu\subseteq\hat s\).
In each clause, finitely many sentences of \(\mu\) are named as
hypotheses. Each of them is a member of \(s\), or \(\varphi\), or a member
of \(\Delta_{\mathbb Q}\), so each mentions only finitely many constants
of \(C\cup D\). Let \(K\) be the finite set of these constants and
\(s_1=s[K]\). By Step~4, \(s_1\in\mathcal S\), and every named hypothesis
lies in \(T(s_1)\). By the truth fact, every named hypothesis holds in the
expansion of every \(\alpha\)-realization of \(s_1\) with
\(\alpha\geq1\). Also \(\mu\subseteq\hat s\subseteq\hat s_1\). So, to find
a member of \(\Sigma\) containing \(\mu\cup\{\chi\}\), it suffices to
find \(s_2\in\mathcal S\) with \(s_1\subseteq s_2\) and \(\chi\in s_2\),
because then \(\mu\cup\{\chi\}\subseteq\hat s_2\).
\begin{enumerate}
\item[C0] A subset of \(\mu\) is a subset of \(\hat s\), so it is in
      \(\Sigma\).
\item[C1] Suppose \(\psi\) and \(\neg\psi\) are both in \(\mu\). Both are
      named hypotheses, and \(s_1\) has a \(1\)-realization, whose
      expansion would satisfy both. This is impossible.
\item[C2] Let \(\neg\psi\in\mu\). The formal negation \({\sim}\psi\)
      has the same constants as \(\psi\), by induction on
      \cref{def:p03-w03-p03-039-formal-negation}, and it holds wherever
      \(\neg\psi\) holds (\cref{thm:p03-w03-p03-039-formal-negation}). By
      the adding fact, \(s_2=s_1\cup\{{\sim}\psi\}\in\mathcal S\).
\item[C3] Let \(\bigwedge_{\psi\in X}\psi\in\mu\) and \(\psi\in X\). The
      constants of \(\psi\) are among those of the conjunction, and
      \(\psi\) holds wherever the conjunction holds. So
      \(s_2=s_1\cup\{\psi\}\in\mathcal S\).
\item[C4] Let \(\bigvee_{\psi\in X}\psi\in\mu\). Then \(X\) is countable.
      For each \(\alpha\) with \(1\leq\alpha<\omega_1\), choose an
      \(\alpha\)-realization \(\mathcal R_\alpha\) of \(s_1\) and a member
      \(\psi_\alpha\in X\) true in its expansion. The sets
      \(\{\alpha:\psi_\alpha=\psi\}\) for \(\psi\in X\) are countably many
      and cover \(\omega_1\setminus\{0\}\). Since \(\omega_1\) is not a
      countable union of countable sets, one of them, say for \(\psi\), is
      uncountable. It is then unbounded in \(\omega_1\), because every
      bounded subset of \(\omega_1\) is countable. For any
      \(1\leq\alpha<\omega_1\), pick \(\alpha'\geq\alpha\) with
      \(\psi_{\alpha'}=\psi\). By monotonicity \(\mathcal R_{\alpha'}\) is
      an \(\alpha\)-realization of \(s_1\) whose expansion satisfies
      \(\psi\). It is therefore an \(\alpha\)-realization of
      \(s_1\cup\{\psi\}\), because the constants of \(\psi\) are among
      those of the disjunction. The case \(\alpha=0\) follows by
      monotonicity. So \(s_2=s_1\cup\{\psi\}\in\mathcal S\) with
      \(\psi\in X\).
\item[C5] Let \(\forall v\,\psi(v)\in\mu\) and \(c\in C\). Let
      \(s_1'=s_1[\{c\}]\in\mathcal S\). In the expansion of every
      \(\alpha\)-realization of \(s_1'\) with \(\alpha\geq1\), the
      universal sentence holds, so \(\psi(c)\) holds. The constants of
      \(\psi(c)\) are those of \(\psi(v)\) and possibly \(c\). So
      \(s_2=s_1'\cup\{\psi(c)\}\in\mathcal S\).
\item[C6] Let \(\exists v\,\psi(v)\in\mu\). Choose
      \(c\in C\setminus C(s_1)\), which exists because \(C\) is infinite
      and \(C(s_1)\) is finite. The constant \(c\) does not occur in
      \(\psi(v)\), whose constants are in \(K\). Given an
      \(\alpha\)-realization of \(s_1\) with \(\alpha\geq1\), its expansion
      satisfies the existential sentence. Choose a witness \(a\) and extend
      \(g\) by \(c\mapsto a\). The sentences of \(s_1\) are unaffected, and
      \(\psi(c)\) holds. This gives an \(\alpha\)-realization of
      \(s_1\cup\{\psi(c)\}\), and the case \(\alpha=0\) follows by
      monotonicity. So \(s_2=s_1\cup\{\psi(c)\}\in\mathcal S\).
\item[C7] Let \(t\) be a closed \(\tau^{+}\)-term and \(c,d\in C\).
      (a) If \(c=d\in\mu\), then \(d=c\) holds wherever \(c=d\) does, so
      \(s_2=s_1\cup\{d=c\}\in\mathcal S\). (b) If \(c=t\in\mu\) and
      \(\psi(t)\in\mu\), then in each relevant expansion \(c\) and \(t\)
      denote the same element, so \(\psi(c)\) holds there, by the
      substitution property of satisfaction for closed terms. The
      constants of \(\psi(c)\) are among those of \(\psi(t)\) together
      with \(c\), so they lie in \(K\), and
      \(s_2=s_1\cup\{\psi(c)\}\in\mathcal S\). (c) The term \(t\) is
      finite, so its set \(K_t\) of constants from \(C\cup D\) is finite.
      Let \(s_1'=s_1[K_t]\in\mathcal S\) and choose
      \(e\in C\setminus C(s_1')\). In an \(\alpha\)-realization of
      \(s_1'\), the term \(t\) denotes an element. Extend \(g\) by sending
      \(e\) to it. Then \(e=t\) holds and \(s_1'\) is unaffected, and the
      case \(\alpha=0\) follows by monotonicity. So
      \(s_2=s_1'\cup\{e=t\}\in\mathcal S\).
\end{enumerate}
So \(\Sigma\) is a nonempty consistency property.

\emph{Step 6: the model.} Let \(T=\{\varphi\}\cup\Delta_{\mathbb Q}\), a
countable set of \(\tau^{+}\)-sentences. If \(\mu\in\Sigma\) with
\(\mu\subseteq\hat s\), then \(\mu\cup\{\chi\}\subseteq\hat s\) for every
\(\chi\in T\), so \(\mu\cup\{\chi\}\in\Sigma\). By
\cref{thm:p03-w03-p03-039-extended-model-existence}, \(T\) has a countable
model \(\mathfrak N^{+}\). Let \(\mathfrak N\) be its \(\tau\)-reduct and
\(h(q)=d_q^{\mathfrak N^{+}}\). Then \(\mathfrak N\entails\varphi\) by the
first fact, its domain is countable, and for \(q<r\) the sentence
\(d_q<d_r\) of \(T\) gives \(h(q)<^{\mathfrak N}h(r)\). With Step~1 this
proves the theorem.
\end{proof}
\end{annexitem}

\begin{annexitem}{annex:p03-w03-p03-041-bounded}{thm:p03-w03-p03-041-bounded}{Single sentences pin down only countable ordinals}
\begin{proof}
\emph{Part 1.} Suppose the conclusion fails. Then for every
\(\alpha<\omega_1\) some model \(\mathfrak M_\alpha\) of \(\varphi\) has
field of order type \(\delta_\alpha\geq\alpha\). Let \(f\) be the
isomorphism of \(\delta_\alpha\) onto the field, ordered by
\(<^{\mathfrak M_\alpha}\), and let \(e\) be its restriction to
\(\alpha\). Then \(e(\xi)<^{\mathfrak M_\alpha}e(\eta)\) whenever
\(\xi<\eta<\alpha\). By \cref{thm:p03-w03-p03-041-undefinability} there
are \(\mathfrak N\entails\varphi\) and
\(h\colon\mathbb Q\to\dom\mathfrak N\) with \(h(q)<^{\mathfrak N}h(r)\)
for \(q<r\). Put \(a_n=h(1/(n+1))\). Then \(a_{n+1}<^{\mathfrak N}a_n\) for
every \(n\), so every \(a_n\) lies in the field of \(<^{\mathfrak N}\). By
hypothesis this field is well-ordered, so the nonempty set
\(\{a_n:n<\omega\}\) has a least element \(a_k\). Then \(a_k=a_{k+1}\) or
\(a_k<^{\mathfrak N}a_{k+1}\). Together with \(a_{k+1}<^{\mathfrak N}a_k\),
the first contradicts irreflexivity and the second asymmetry of a strict
linear order. So there is \(\alpha<\omega_1\) such that every model has
field type less than \(\alpha\).

\emph{Part 2.} The lower bound \(w(L_{\omega_1\omega})\geq\omega_1\) is
\cref{thm:p03-w03-p03-041-countable-ordinals}, and
\(w_{\aleph_0}\geq w_1\) because a set of at most one sentence is
countable. It remains to show
\(w_{\aleph_0}(L_{\omega_1\omega})\leq\omega_1\). Let \(\Phi\) be a
countable set of \(L_{\omega_1\omega}\)-sentences of a signature \(\tau\)
that pins down \(\beta\) via \(<\). If \(\Phi\) were empty, condition~(i)
of \cref{def:p03-w03-p03-041-pinning} would fail: the \(\tau\)-structure
with domain \(\omega\), with \(<\) interpreted as the reverse of the usual
order and every other symbol interpreted arbitrarily, has a field with no
least element. So \(\Phi\) is nonempty, and its conjunction
\(\varphi=\bigwedge\Phi\) is an \(L_{\omega_1\omega}\)-sentence with the
same models as \(\Phi\). By condition~(i), \(\varphi\) satisfies the
hypothesis of part~1, which gives some \(\alpha<\omega_1\). By
condition~(ii) some model has field of type \(\beta\), so
\(\beta<\alpha<\omega_1\). Every ordinal pinned down by a countable set of
sentences is therefore countable, and the supremum of these ordinals is at
most \(\omega_1\). With the lower bound,
\(w(L_{\omega_1\omega})=w_{\aleph_0}(L_{\omega_1\omega})=\omega_1\).
\end{proof}
\end{annexitem}
